\section{Discussion}
\label{sec:discussion}

\subsection{When a submission can be refuted, and when it cannot}

Independent checking of national greenhouse-gas submissions usually ends inconclusively. Two
estimates differ; both are models; neither is a measurement; and the comparison records a
disagreement without resolving it. This inventory's fugitives sector is a clean example, where
a factor of \num{1.3} against EDGAR settles nothing because both estimates choose within the
same published IPCC range and neither has observed a Pakistani gas system.

Railways is the opposite case, and the difference is worth stating precisely because it is not
about rigour. It is about what the category rests on.

A fugitive emission factor is a statement about infrastructure nobody has measured. A railway
fuel figure is a quantity a state enterprise metered, paid for and published, and the carbon
in it is fixed by chemistry. When a submitted figure requires ten times that fuel, the
disagreement is not between two models---it is between a figure and an audited quantity. And
when a second, independent statistic in the same publication puts the implied fuel intensity at
\SI{46}{\litre} per train-kilometre, the refutation no longer depends on trusting the fuel
account either.

The general form: a category can be refuted from outside when its activity datum is a metered
physical quantity and a second published statistic bounds the same quantity independently.
Railways has both. Most categories have neither, which is why this kind of result is rare
rather than because verification is usually done badly.

\subsection{What probably went wrong in the submission}

This paper does not know how the BUR's \SI{4.06}{\mega\tonne} was produced, and
speculating about it is not the same as refuting it, so that refutation stands on the
arithmetic alone. But the current CRT localises its own version of the error, because a
CRT publishes intermediates the BUR did not: the activity cell is \num{12.6} times the
metered fuel---a magnitude characteristic of a unit error, litres read as something
larger---and the implied-factor cells show the road table where the rail table belongs.
Two errors, separable in the filing's own arithmetic, each of which the operator's two
published statistics would have exposed immediately.

The constructive point is procedural rather than diagnostic. Had the submission published its
activity datum alongside its emissions, the disagreement would have been visible immediately
and resolvable by anyone. It is the absence of the intermediate quantity, not the estimate,
that makes the figure hard to check---and publishing implied fuel per train-kilometre alongside
the emissions would cost nothing and make the category self-auditing.

\subsection{What global inventories are missing, demonstrated}

The ELK comparison is worth generalising. A careful global fleet model, given no Pakistani
activity data, lands a factor of two low nationally and inverts the passenger--freight
composition entirely---five grid cells of passenger rail for an operator that runs four
passenger train-kilometres for every freight one. The two numbers that would have fixed it,
annual fuel and train-kilometres, have been published in the operator's Year Book every
year since the 1950s. The marginal value of nationally reported operator data to global
inventories is usually asserted; here it is measurable, in both magnitude and composition.

\subsection{The third instance of the parallel-table trap}

Volume~2 Chapter~3 contains several near-duplicate default tables. Their \COtwo{} rows agree,
because \COtwo{} depends only on the carbon in the fuel; their \CHfour{} and \NtwoO{} rows do
not, because those depend on the engine and the duty cycle.

This inventory has now hit that structure three times: residential fuel against the power
table, power against several, and rail against road here. In every case the failure would be
silent---the dominant gas is unaffected, the fuel reconciliation still closes, and the error
appears only in a minor gas at a magnitude that looks like ordinary uncertainty.

The defence that generalises is not vigilance. It is to publish the counterfactual. This sector
computes what it would report on the wrong table, \SI{0.363}{\mega\tonne} against
\SI{0.394}{}, and writes it into its own output as a labelled diagnostic row. A reader can then
see the size of the choice without re-deriving it, and a future edit that silently switched
tables would move a number that is checked rather than one that is not.

\subsection{Small categories and the attributions they settle}

At \SI{0.09}{\percent} of national emissions, railways would fail any threshold test for
inclusion. Its value to the inventory is not its own magnitude but the attribution it
settles: the national high-speed-diesel figure covers road and rail together, and before the
road sector re-anchored on the petroleum industry's customer classes, the road category
absorbed the whole of it. The two fuels are now separate rows at the source, and this
sector's metered account is the independent check that the Railways row is right.

That is a general property of small transport categories. They are frequently the reason a
large category's activity data can be stated correctly, and dropping them on grounds of size
leaves the large category quietly wrong by the amount that was dropped. Rail was worth about
a third of a megatonne to road transport, which is roughly the entire railway sector---the
correction to the neighbour was the same size as the category making it, and it now lives in
the source data rather than running as a deduction.
